\documentclass[11pt,a4paper]{article}

% ============================================================
% PAQUETES
% ============================================================
\usepackage[utf8]{inputenc}
\usepackage[T1]{fontenc}
\usepackage[spanish,es-nodecimaldot,es-noquoting]{babel}
\usepackage{lmodern}
\usepackage{geometry}
\usepackage{amsmath,amssymb}
\usepackage{booktabs}
\usepackage{tabularx}
\usepackage{array}
\usepackage{enumitem}
\usepackage{ragged2e}
\usepackage{xcolor}
\usepackage{tcolorbox}
\tcbuselibrary{skins,breakable}
\usepackage{fancyhdr}
\usepackage{lastpage}
\usepackage{microtype}
\usepackage{hyperref}

% ============================================================
% CONFIGURACIÓN DE PÁGINA
% ============================================================
\geometry{
  top=2.0cm,
  bottom=2.0cm,
  left=2.0cm,
  right=2.0cm
}

\setlength{\parindent}{0pt}
\setlength{\parskip}{5pt}
\renewcommand{\arraystretch}{1.2}

% ============================================================
% COLORES INSTITUCIONALES Y TÉCNICOS
% ============================================================
\definecolor{unlgreen}{RGB}{0,91,65}
\definecolor{azulOscuro}{RGB}{12,43,92}
\definecolor{azulMedio}{RGB}{30,80,140}
\definecolor{verdeBosque}{RGB}{27,107,74}
\definecolor{acentoDorado}{RGB}{197,143,39}
\definecolor{rojoAlerta}{RGB}{165,40,40}
\definecolor{grisFondo}{RGB}{247,249,252}
\definecolor{grisBorde}{RGB}{209,217,230}
\definecolor{grisTexto}{RGB}{55,65,75}

% ============================================================
% CONFIGURACIÓN DE HYPERREF
% ============================================================
\hypersetup{
  colorlinks=true,
  linkcolor=azulOscuro,
  urlcolor=azulMedio,
  citecolor=verdeBosque
}

% ============================================================
% ENCABEZADO Y PIE DE PÁGINA
% ============================================================
\pagestyle{fancy}
\fancyhf{}
\setlength{\headheight}{22pt}
\renewcommand{\headrulewidth}{0.8pt}
\renewcommand{\headrule}{\hbox to\headwidth{\color{unlgreen}\leaders\hrule height \headrulewidth\hfill}}
\fancyhead[L]{\small\textbf{\color{unlgreen}UNIVERSIDAD NACIONAL DE LOJA} $\vert$ FARNR}
\fancyhead[R]{\small\textbf{Estadística Descriptiva} $\vert$ Guía 1: Bloque C}
\fancyfoot[L]{\footnotesize Carrera de Ingeniería Ambiental $\vert$ Guillermo Chuncho}
\fancyfoot[R]{\footnotesize Pág. \thepage\ de \pageref{LastPage}}

\fancypagestyle{plain}{%
  \fancyhf{}%
  \renewcommand{\headrulewidth}{0pt}%
  \fancyfoot[L]{\footnotesize Carrera de Ingeniería Ambiental $\vert$ Guillermo Chuncho}%
  \fancyfoot[R]{\footnotesize Pág. \thepage\ de \pageref{LastPage}}%
}

% ============================================================
% COLUMNAS TABULARES
% ============================================================
\newcolumntype{Y}{>{\RaggedRight\arraybackslash}X}
\newcolumntype{C}[1]{>{\centering\arraybackslash}p{#1}}
\newcolumntype{L}[1]{>{\RaggedRight\arraybackslash}p{#1}}

% ============================================================
% CAJAS DE DISEÑO
% ============================================================
\newtcolorbox{bannerbox}[1][]{
  colback=azulOscuro, colframe=azulOscuro, arc=4pt, boxrule=1pt,
  fonttitle=\bfseries\color{white}, colbacktitle=azulOscuro,
  top=6pt, bottom=6pt, left=8pt, right=8pt, breakable, #1
}

\newtcolorbox{objetivobox}[1][]{
  colback=green!3!white, colframe=unlgreen, arc=3pt, boxrule=1pt,
  title=\textbf{Propósito y Resultados de Aprendizaje},
  fonttitle=\bfseries\color{white}, colbacktitle=unlgreen,
  top=5pt, bottom=5pt, left=8pt, right=8pt, breakable, #1
}

\newtcolorbox{teoriabox}[1][]{
  colback=grisFondo, colframe=azulMedio, arc=3pt, boxrule=1pt,
  title=\textbf{Fundamento Conceptual -- Bloque C: Diseños de Muestreo Ambiental},
  fonttitle=\bfseries\color{white}, colbacktitle=azulMedio,
  top=5pt, bottom=5pt, left=8pt, right=8pt, breakable, #1
}

\newtcolorbox{casobox}[2][]{
  colback=white, colframe=azulOscuro, arc=3pt, boxrule=1pt,
  title=\textbf{#2}, fonttitle=\bfseries\color{white}, colbacktitle=azulOscuro,
  top=6pt, bottom=6pt, left=8pt, right=8pt, breakable, #1
}

\newtcolorbox{alertbox}[1][]{
  colback=red!3!white, colframe=rojoAlerta, arc=3pt, boxrule=1pt,
  title=\textbf{Advertencia Metodológica: Error Muestral vs. Sesgo Sistemático},
  fonttitle=\bfseries\color{white}, colbacktitle=rojoAlerta,
  top=5pt, bottom=5pt, left=8pt, right=8pt, breakable, #1
}

\begin{document}

% ============================================================
% ENCABEZADO INSTITUCIONAL
% ============================================================
\begin{bannerbox}
\centering
{\large\bfseries\color{white} UNIVERSIDAD NACIONAL DE LOJA}\\[2pt]
{\small\bfseries\color{white} FACULTAD AGROPECUARIA Y DE RECURSOS NATURALES RENOVABLES}\\[2pt]
{\footnotesize\color{white} CARRERA DE INGENIERÍA AMBIENTAL $\cdot$ ASIGNATURA: ESTADÍSTICA DESCRIPTIVA}\\[4pt]
\rule{\linewidth}{0.5pt}\\[4pt]
{\large\bfseries\color{white} GUÍA DE ESTUDIO 1: ELEMENTOS DE MUESTREO Y DISEÑO DEL ESTUDIO}\\[2pt]
{\footnotesize\color{white} BLOQUE C: POBLACIÓN, MUESTRA, MARCO MUESTRAL, PARÁMETROS, SESGO Y MÉTODOS DE MUESTREO}\\[2pt]
{\scriptsize\color{white} Docente: Guillermo Chuncho $\cdot$ Contacto: \texttt{carguille2@gmail.com} $\cdot$ Ciclo 2026--2027}
\end{bannerbox}

\vspace{0.3cm}

\begin{objetivobox}
\begin{itemize}[leftmargin=*,noitemsep]
  \item \textbf{Resultado de Aprendizaje:} Dominar la transición epistemológica y técnica que va desde la pregunta de investigación ambiental hasta la estructuración de la matriz de datos muestrales representativa.
  \item \textbf{Competencias Específicas:} Diferenciar con rigor técnico población objetivo, población accesible, marco muestral, unidad muestral y unidad de análisis. Distinguir formalmente parámetros de estadísticos, cuantificar el error muestral frente al sesgo sistemático, y seleccionar el método de muestreo adecuado (probabilístico vs. no probabilístico) para la cuenca del cantón Loja.
  \item \textbf{Recurso Interactivo Web:} Explore la simulación Monte Carlo de parámetros vs. estadísticos en la aplicación oficial: \url{https://informacion-estadistica-unl.streamlit.app}.
\end{itemize}
\end{objetivobox}

\vspace{0.2cm}

% ============================================================
% SECCIÓN 1: ELEMENTOS OPERATIVOS DEL MUESTREO
% ============================================================
\section*{1. Los Cinco Componentes Operativos del Muestreo}

En investigaciones ambientales, uno de los errores más frecuentes consiste en confundir el ámbito conceptual de estudio con la unidad física de medición o con el registro operativo de selección.

\begin{table}[h!]
\centering
\small
\begin{tabularx}{\linewidth}{p{3.2cm}p{4.2cm}p{4.2cm}Y}
\toprule
\textbf{Elemento} & \textbf{Definición Técnica} & \textbf{Pregunta Clave} & \textbf{Ejemplo Cantón Loja} \\
\midrule
\textbf{Población Objetivo} ($N$) & Totalidad de elementos o eventos sobre los cuales se desea formular inferencias. & ¿A qué universo completo se desea extrapolar? & Todos los tramos fluviales continuos del Río Malacatos en el cantón Loja en estiaje. \\
\midrule
\textbf{Población Accesible} & Subconjunto de la población objetivo que es física y operativamente muestreable. & ¿Qué elementos están realmente disponibles? & 140 tramos fluviales con servidumbre de paso público o puentes comunitarios. \\
\midrule
\textbf{Marco Muestral} & Listado formal, base cartográfica o catálogo operativo desde donde se extrae la muestra. & ¿Cuál es la lista tangible y numerada de sorteo? & Archivo shapefile SIG con 140 tramos elegibles codificados de \texttt{ID\_001} a \texttt{ID\_140}. \\
\midrule
\textbf{Unidad Muestral} & Entidad operativa seleccionada en el sorteo o procedimiento de selección. & ¿Qué objeto o punto se extrae en cada sorteo? & El tramo fluvial de 50 metros georreferenciado en el sorteo. \\
\midrule
\textbf{Unidad de Análisis} & Objeto indivisible donde se mide la variable o parámetro ambiental. & ¿En qué ente se mide la variable en laboratorio? & La alícuota de 500 mL de agua sometida a análisis de turbidez y DBO$_5$. \\
\bottomrule
\end{tabularx}
\end{table}

\begin{alertbox}
\textbf{El Defecto de Subcobertura en el Marco Muestral:}\\
Si el marco muestral no coincide exactamente con la población objetivo, se produce un \textbf{sesgo de cobertura (subregistro)}. Por ejemplo, si se evalúa la contaminación en huertos periurbanos utilizando únicamente el catastro formal de predios con título, se excluyen sistemáticamente los asentamientos informales, donde la vulnerabilidad y la exposición a metales pesados suelen ser mayores. Aumentar el tamaño de la muestra jamás corrige este defecto.
\end{alertbox}

\vspace{0.2cm}

% ============================================================
% SECCIÓN 2: PARÁMETRO VS ESTADÍSTICO
% ============================================================
\section*{2. Parámetro Poblacional vs. Estadístico Muestral}

La inferencia estadística descansa sobre la dualidad entre la \textbf{verdad global inobservable} y la \textbf{evidencia muestral calculada}:

\begin{table}[h!]
\centering
\small
\begin{tabularx}{\linewidth}{p{3.5cm}p{4.2cm}p{4.2cm}Y}
\toprule
\textbf{Criterio} & \textbf{Parámetro Poblacional ($\theta$)} & \textbf{Estadístico Muestral ($\hat{\theta}$)} & \textbf{Relación Inferencial} \\
\midrule
\textbf{Ámbito de Origen} & Proviene de la \textbf{Población total ($N$)}. & Proviene de la \textbf{Muestra observada ($n$)}. & $\hat{\theta}$ estima a $\theta$ con incertidumbre. \\
\midrule
\textbf{Naturaleza Matemática} & Constante fija. Prácticamente siempre desconocida (requiere censo total). & Variable aleatoria. Cambia de una muestra a otra según la ley de probabilidad. & Su varianza muestral es el error estándar: $\sigma_{\hat{\theta}}$. \\
\midrule
\textbf{Simbología Habitual} & Letras griegas: $\mu$ (media), $\sigma^2$ (varianza), $\sigma$ (desviación), $p$ (proporción). & Letras latinas: $\bar{x}$ (media), $s^2$ (varianza), $s$ (desviación), $\hat{p}$ (proporción). & Notación formal obligatoria en artículos y reportes técnicos. \\
\midrule
\textbf{Ejemplo Numérico} & $\mu = 45.00$ NTU (turbidez real desconocida de toda la cuenca). & $\bar{x} = 43.80$ NTU (calculado en $n = 30$ estaciones de campo). & Error muestral puntual: $\vert \bar{x} - \mu \vert = 1.20$ NTU. \\
\bottomrule
\end{tabularx}
\end{table}

\vspace{0.2cm}

% ============================================================
% SECCIÓN 3: ERROR MUESTRAL VS SESGO SISTEMÁTICO
% ============================================================
\section*{3. Error Muestral vs. Sesgo Sistemático (Bias)}

\begin{figure}[h!]
\centering
\small
\begin{tabularx}{\linewidth}{YY}
\begin{tcolorbox}[colback=blue!3!white,colframe=azulMedio,title=\textbf{Error Muestral ($\epsilon = \hat{\theta} - \theta$)}]
\begin{itemize}[leftmargin=*,noitemsep]
  \item \textbf{Origen:} Fluctuación propia del azar por observar una fracción y no toda la población.
  \item \textbf{Propiedad:} Es aleatorio y centrado en cero: $\mathbb{E}[\hat{\theta} - \theta] = 0$.
  \item \textbf{Efecto de $n$:} \textbf{Disminuye} conforme $n$ aumenta, a razón de $1/\sqrt{n}$.
  \item \textbf{Impacto:} Afecta la \textbf{precisión} (dispersión de estimaciones).
\end{itemize}
\end{tcolorbox}
&
\begin{tcolorbox}[colback=red!3!white,colframe=rojoAlerta,title=\textbf{Sesgo Sistemático ($Bias = \mathbb{E}[\hat{\theta}] - \theta$)}]
\begin{itemize}[leftmargin=*,noitemsep]
  \item \textbf{Origen:} Defectos estructurales en calibración de sensores, diseño muestral o marco.
  \item \textbf{Propiedad:} Es unidireccional y persistente: $\mathbb{E}[\hat{\theta}] \neq \theta$.
  \item \textbf{Efecto de $n$:} \textbf{NO DISMINUYE}. Con $n = 500.000$ el sesgo se perpetúa intacto.
  \item \textbf{Impacto:} Destruye la \textbf{exactitud y validez científica} del estudio.
\end{itemize}
\end{tcolorbox}
\end{tabularx}
\end{figure}

\vspace{0.2cm}

% ============================================================
% SECCIÓN 4: MÉTODOS DE MUESTREO AMBIENTAL
% ============================================================
\section*{4. Clasificación Rigurosa de Métodos de Muestreo}

\subsection*{4.1 Diseños Probabilísticos (Permiten Inferencia y Cálculo de Confianza)}
\begin{enumerate}[leftmargin=*]
  \item \textbf{Aleatorio Simple (MAS):} Cada unidad tiene probabilidad idéntica ($1/N$). Requiere marco muestral completo y variabilidad uniforme. \textit{Ejemplo: Sorteo de 35 parcelas de una cuadrícula de 400 predios.}
  \item \textbf{Estratificado:} La población se divide en estratos internamente homogéneos y heterogéneos entre sí.
  \begin{itemize}
    \item \textit{Proporcional:} $n_h = n \cdot \frac{N_h}{N}$. Mantiene la representatividad poblacional relativa.
    \item \textit{No Proporcional / Óptima:} Se asigna mayor cuota a estratos raros pero de alto valor ecológico (ej. turberas del Parque Podocarpus) o de mayor variabilidad interna ($\sigma_h$).
  \end{itemize}
  \item \textbf{Sistemático:} Se elige un inicio aleatorio $r \in [1, k]$ y luego cada $k = N/n$ elementos continuos.
  \textit{Riesgo crítico:} Si $k$ coincide con la periodicidad del fenómeno (ej. vertidos industriales semanales), se genera un sesgo masivo.
  \item \textbf{Por Conglomerados (Clusters):} Se seleccionan grupos completos (microcuencas, parroquias) y se censan todas sus unidades internas para minimizar costos de transporte.
  \item \textbf{Polietápico:} Combinación jerárquica de sorteos probabilísticos en fases sucesivas (ej. Parroquias $\to$ Barrios $\to$ Intersecciones viales).
\end{enumerate}

\subsection*{4.2 Diseños No Probabilísticos (Exploratorios o de Juicio Técnico)}
\begin{enumerate}[leftmargin=*]
  \item \textbf{Por Conveniencia:} Selección por proximidad o bajo costo (ej. tomar muestras de agua solo donde hay puentes vehiculares). \textit{Alto riesgo de sesgo.}
  \item \textbf{Por Cuotas:} Se fijan cupos por categoría pero la selección final no es aleatoria.
  \item \textbf{Intencional o por Juicio Experto:} El investigador escoge sitios críticos conocidos (ej. inspeccionar las 3 descargas mineras más conflictivas).
  \item \textbf{Bola de Nieve (Snowball):} Participantes iniciales contactan y refieren a otros miembros de poblaciones clandestinas o sin registro formal (ej. tala no autorizada de madera nativa o quemas agrícolas ilegales).
\end{enumerate}

\vspace{0.2cm}

% ============================================================
% SECCIÓN 5: LOS CUATRO CASOS DE ESTUDIO DE LOJA
% ============================================================
\section*{5. Análisis Metodológico de los Casos de Estudio en Loja}

\begin{casobox}{Caso 1: Calidad de Agua en la Microcuenca del Río Malacatos}
\textbf{Objetivo:} Evaluar coliformes fecales y DBO$_5$ en tramos periurbanos y rurales durante el invierno.\\
\textbf{Población Objetivo:} Todos los tramos del Río Malacatos en el cantón Loja ($N$ continuo).\\
\textbf{Población Accesible:} 140 tramos con senderos públicos accesibles y seguros.\\
\textbf{Marco Muestral:} Mapa SIG con 140 tramos identificados (\texttt{ID\_001} a \texttt{ID\_140}).\\
\textbf{Unidad Muestral:} Tramo de 50 metros seleccionado en el sorteo. $\vert$ \textbf{Unidad de Análisis:} Muestra de 500 mL analizada en laboratorio.\\
\textbf{Diseño Recomendado:} \textbf{Estratificado por Piso Altitudinal} (alto rural, medio urbano, bajo periurbano).
\end{casobox}

\begin{casobox}{Caso 2: Generación y Separación de Residuos Domiciliarios}
\textbf{Objetivo:} Estimar la proporción cantonal ($p$) de hogares que realizan separación en origen.\\
\textbf{Población Objetivo:} Todos los hogares habitados de las áreas urbanas de Loja ($N \approx 55.000$).\\
\textbf{Marco Muestral:} Catastro de contribuyentes de aseo y agua potable.\\
\textbf{Unidad Muestral:} Vivienda o predio sorteado. $\vert$ \textbf{Unidad de Análisis:} El hogar habitante (kg/persona/día).\\
\textbf{Diseño Recomendado:} \textbf{Estratificado Proporcional} por circuito barrial o estrato socioeconómico.
\end{casobox}

\begin{casobox}{Caso 3: Ocurrencia de Incendios Forestales (Datos Satelitales VIIRS 375 m)}
\textbf{Objetivo:} Estimar la proporción cantonal de superficie afectada por incendios en el estiaje.\\
\textbf{Población Objetivo:} 7.997 celdas espaciales de $500 \times 500$ m del cantón Loja.\\
\textbf{Marco Muestral:} Malla raster/vectorial con 7.997 filas identificadas de \texttt{cell\_id = 1} a \texttt{7997}.\\
\textbf{Unidad Muestral y de Análisis:} La celda de $500 \times 500$ m (variable binaria 0/1 y severidad dNBR).\\
\textbf{Diseño Recomendado:} \textbf{Estratificado por Tipo de Cobertura Vegetal} (páramo, matorral, pino, pastos).
\end{casobox}

\begin{casobox}{Caso 4: Monitoreo de Material Particulado PM$_{2.5}$ en el Anillo Urbano Céntrico}
\textbf{Objetivo:} Determinar el promedio anual de PM$_{2.5}$ en horas pico de flujo vehicular.\\
\textbf{Población Objetivo:} Todas las horas diurnas del año en la malla vial céntrica ($N = 5.110$ horas).\\
\textbf{Marco Muestral:} Calendario cronológico de días y puntos georreferenciados para instalación.\\
\textbf{Unidad Muestral:} Tramo vial y franja horaria de 1 hora. $\vert$ \textbf{Unidad de Análisis:} Metro cúbico de aire aspirado.\\
\textbf{Diseño Recomendado:} \textbf{Sistemático Temporal} (cada $k = 3$ días en horarios punta fijos).
\end{casobox}

\end{document}
